\section{Introduction}
\label{sec:introduction}

A paper reports that reinforcement learning post-training moved a benchmark by two points. Whether
that is a result depends on a number nobody has: how much the same recipe would have moved it with
a different seed. The problem is documented \citep{henderson2018matters, hochlehnert2025sober} and
the only instrument the field has is replication, which at current RL compute is close to
unaffordable. So the number usually goes unmeasured.

Start from the obvious model. Gradient noise perturbs each update, perturbations accumulate, two
runs execute a random walk away from each other, and reproducibility decays with run length. That
picture is implicit in how run-to-run variance is discussed, and it is what a noisy gradient
suggests.

It is wrong, and the way it is wrong is the subject of this paper. Noise from one update does not
survive to the end of training unchanged. It is passed through every update that follows, and where
the shared objective is locally concave those updates shrink it. What reaches the end is a
convolution of injected noise with the sensitivity of the remaining trajectory, and that object is
computable.

\paragraph{Contributions.}
\begin{itemize}
  \item \textbf{A finite-time covariance law} (Theorem~\ref{thm:propagation}).
  $S_{t+1} = A_t S_t A_t^{\!\top} + Q_t$, giving expected divergence between seeds as
  $\tr(F_T S_T)$ and attributing it to individual updates through a per-update kernel. In an
  exactly solvable policy it predicts measured divergence to $9\%$ at the median over twelve
  settings, against $2.0\times$ for accumulation and two orders of magnitude for a single-timescale
  approximation.

  \item \textbf{The contraction is caused by learning.} Replacing rewards with coin flips, which
  leaves noise and batch composition untouched, turns a flat divergence trace into one that grows
  $10.7\times$ over eighty updates. This separates the result from a claim that seeds merely end up
  nearby.

  \item \textbf{The machinery that makes the injected noise measurable}: an exact finite-$G$ mean
  and covariance for the count-based advantage estimators under binary rewards
  (Propositions~\ref{prop:mean} and \ref{prop:cov}), a decomposition of the gradient noise into
  between-prompt and within-prompt terms, and an estimator that recovers both from the microbatch
  gradients a trainer already accumulates. The same decomposition fixes the compute-optimal group
  size (Theorem~\ref{thm:split}), which we validate prospectively.

  \item \textbf{An explicit account of what does not hold.} A single contraction timescale fails by
  two orders of magnitude, and we report it as the comparator rather than the law. A predicted
  square-root relation between policy divergence and outcome spread is not supported by our data
  and is withdrawn. Drift-targeted control, which Section~\ref{sec:drift} otherwise recommends,
  destroys runs at saturation. The largest model we measure is $0.5$B, which bounds what any of
  this establishes about frontier-scale RLVR.
\end{itemize}

Every experiment here runs on one laptop, which sets the ceiling on the last point and is stated
plainly in Section~\ref{sec:discussion} rather than left for a reader to infer.
